% !TEX TS-program = lualatex+se
% !TEX encoding = UTF-8 Unicode

\documentclass[17pt, twoside, a4paper]{extreport}
% -----------------------------------------------------------------------------------
% ------------------------- USTAWIENIA STRONY--------------------------
% -----------------------------------------------------------------------------------
\usepackage{geometry}
\geometry{
	left=10mm,
	right=15mm,
	top=15mm,
	bottom=15mm,
	marginparwidth=0mm,
	marginparsep=0mm
	}
\newcommand{\wstawplik}[1]{\input{../../include/#1}}%skraca kod
\wstawplik{headertest}

\begin{document}
% -----------------------------------------------------------------------------------
% ------------------------- STRONA TYTUŁOWA  --------------------------
\thispagestyle{empty}
\begin{center}
	\includegraphics[width=1.0\linewidth]{HolyThursday03.jpg}
	\vfill
	\podtytulfont{Ad Matutinum et Laudes}
	\tytulfont{ Feria Quinta\\in C\oe{}na Domini}
	\vfill
	\large\textit{Liber usualis, Solesmes, 1961, p.626}
\end{center}%
\newpage
\thispagestyle{empty}
\wstawplik{modlitwy_wstepne_has}
\newpage
\setcounter{page}{1}
\fancyhead{}% wyczyść wszystkie pola nagłówka
\pagestyle{plain}% nagłówek i stopka pusta - UWAGA: przedefiniowany styl przez polecenie fancypagestyle w header.tex	
%\setlength{\parindent}{0pt}
\gresetbracerendering{font}
\sekcja{In I. Nocturno}{}

%% -------- 1. Ant. Zelus -----------
\gresetinitiallines{1}\greannotation{1. Ant.}\greannotation{8 c}\gabcsnippet{(c4)ZE(j)lus(j) do(j)mus(j) tu(j_k)æ(ji/jkj') <c><v>*</v></c> com(h)é(j')dit(i) me,(g.) (;) et(f) op(h)pró(j')bri(i)a(g_[uh:l]h) ex(f)pro(h)brán(j')ti(k)um(i') ti(j)bi(h.) (,) ce(h)ci(g)dé(hi)runt(h') su(g)per(fg~) me.(g.) (::)}
\tlumaczR{Gryzła mnie żarliwość~\red{*} domu twego i urągania urągających tobie na mnie spadły.}
\subsekcja{Psalm 68}
\gresetinitiallines{0}\gabcsnippet{(c4)<c><v>\footnotesize{1. }</v></c>Sal(g)vum(h) me(j) fac,(j) <b>De</b>(k jr)us:(j.) <c><v>*</v></c> (:) quón(j)i(j)am(j) in(j)tra(j)vé(j)runt(j) a(j)quæ(j) us(j)que(j) ad(j) á(j)<i><ul>ni</i></ul>(h)<i><ul>mam</i></ul>(j) <b>me</b>(k jr)am.(j.) (::) <c><i>Flex:</i></c>()  tem(j)pés(j)tas(j) a(j)quæ, <c><v>†</v></c>(h. j j j  ::)}
\tlumaczR{Wybaw mię, Boże:~\red{*} boć weszły wody aż do duszy mojéj.}\linespace
\columnratio{0.7}
\setlength{\columnseprule}{0.4pt}
\setlength{\columnsep}{.02\textwidth}
\BeginParallel
\latN{\mred{2.}~Infíxus sum in limo pro\textbf{fún}di:~\red{*} et non \underline{\textit{est}} \underline{\textit{sub}}\textbf{stán}tia.}\polPs{Ulgnąłem w błocie głębokości:~\red{*} i dna niemasz.}%

\latN{\mred{3.}~Veni in altitúdinem \textbf{ma}ris:~\red{*} et tempés \underline{\textit{tas}} \underline{\textit{de}}\textbf{mér}sit me.}\polPs{Przyszedłem na głębokość morską:~\red{*} a nawałność mię ponurzyła.}%

\latN{\mred{4.}~Laborávi clamans, raucæ factæ sunt fauces \textbf{me}æ:~\red{*} defecérunt óculi mei, dum spero in \underline{\textit{De}}\underline{\textit{um}} \textbf{me}um.}\polPs{Spracowałem się wołając, ochrapiało gardło moje:~\red{*} ustały oczy moje, gdy mam nadzieję w Bogu moim.}%

\latN{\mred{5.}~Multiplicáti sunt super capíllos cápitis \textbf{me}i,~\red{*} qui odé\underline{\textit{runt}} \underline{\textit{me}} \textbf{gra}tis.}\polPs{Rozmnożyli się nad włosy głowy mojéj,~\red{*} którzy mię mają w nienawiści bez przyczyny.}%

\latN{\mred{6.}~Confortáti sunt qui persecúti sunt me inimíci mei in\textbf{jús}te:~\red{*} quæ non rápui, tunc \underline{\textit{ex}}\underline{\textit{sol}}\textbf{vé}bam.}\polPs{Zmocnili się, którzy mię niesprawiedliwie prześladowali, nieprzyjaciele moi:~\red{*} czegom nie wydarł, tedym płacił.}%

\latN{\mred{7.}~Deus, tu scis insipiéntiam \textbf{me}am:~\red{*} et delícta mea a te non \underline{\textit{sunt}} \underline{\textit{abs}}\textbf{cón}dita.}\polPs{Boże, ty znasz głupstwo moje:~\red{*} a występki moje nie są tobie tajne.}%

\latN{\mred{8.}~Non erubéscant in me qui exspéctant te, \textbf{Dó}mine,~\red{*} Dómi\underline{\textit{ne}} \underline{\textit{vir}}\textbf{tú}tum}\polPs{Niech nie będą zawstydzeni dla mnie ci, którzy cię oczekiwają, Panie,~\red{*} Panie zastępów.}%

\latN{\mred{9.}~Non confundántur \textbf{su}per me~\red{*} qui quærunt te, \underline{\textit{De}}\underline{\textit{us}} \textbf{Is}raël.}\polPs{Niech nie będą pohańbieni dla mnie,~\red{*} którzy cię szukają, Boże Izraelski.}%
\EndParallel
%\latN{\mred{10.}~Quóniam propter te sustínui op\textbf{pró}brium:~* opéruit confúsio fá\underline{\textit{ci}}\underline{\textit{em}} \textbf{me}am.}\polPs{Bom dla ciebie znaszał urąganie:~\red{*} zelżywość okrywała oblicze moje.}%

\gresetinitiallines{0}\gabcsnippet{(c4)<c><v>\footnotesize{10. }</v></c>Quón(j)i(j)am(j) prop(j)ter(j) te(j) sus(j)tí(j)nu(j)i(j) op(j)<b>pró</b>(k)bri(jr)um:(j.) <c><v>*</v></c>(:) o(j)pé(j)ru(j)it(j) con(j)fú(j)si(j)o(j) fá(j)<i><ul>ci</i></ul>(h)<i><ul>em</i></ul>(j) <b>me</b>(k jr)am.(j.) (::)}%
\vspace{-8mm}
\tlumaczR{Bom dla ciebie znaszał urąganie:~\red{*} zelżywość okrywała oblicze moje.}\linespace
\BeginParallel
\latN{\mred{11.}~Extráneus factus sum frátribus \textbf{me}is,~\red{*} et peregrínus fíliis \underline{\textit{ma}}\underline{\textit{tris}} \textbf{me}æ.}\polPs{Stałem się obcym braciéj mojéj,~\red{*} i cudzoziemcem synom matki mojéj.}%

\latN{\mred{12.}~Quóniam zelus domus tuæ com\textbf{é}dit me:~\red{*} et oppróbria exprobrántium tibi ceci\underline{\textit{dé}}\underline{\textit{runt}} \textbf{su}per me.}\polPs{Bo mię zawisna miłość domu twego gryzła:~\red{*} a urągania urągających tobie spadły na mię.}%

\latN{\mred{13.}~Et opérui in jejúnio ánimam \textbf{me}am:~\red{*} et factum est in oppró\underline{\textit{bri}}\underline{\textit{um}} \textbf{mi}hi.}\polPs{I okryłem w poście duszę moję:~\red{*} i stało mi się urąganiem.}%

\latN{\mred{14.}~Et pósui vestiméntum meum ci\textbf{lí}cium:~\red{*} et factus sum illis \underline{\textit{in}} \underline{\textit{pa}}\textbf{rá}bolam.}\polPs{I oblokłem miasto szaty włosień:~\red{*} i stałem się im przypowieścią.}%

\latN{\mred{15.}~Advérsum me loquebántur, qui sedébant in \textbf{por}ta:~\red{*} et in me psallébant qui bi\underline{\textit{bé}}\underline{\textit{bant}} \textbf{vi}num.}\polPs{Mówili przeciwko mnie, którzy siedzieli w bramie:~\red{*} i śpiewali przeciw mnie, którzy pili wino.}%

\latN{\mred{16.}~Ego vero oratiónem meam ad te, \textbf{Dó}mine:~\red{*} tempus beneplá\underline{\textit{ci}}\underline{\textit{ti}}, \textbf{De}us.}\polPs{Ale ja modlitwę moję do ciebie, Panie:~\red{*} czas upodobania, Boże.}%

\latN{\mred{17.}~In multitúdine misericórdiæ tuæ ex\textbf{áu}di me,~\red{*} in veritáte sa\underline{\textit{lú}}\underline{\textit{tis}} \textbf{tu}æ:}\polPs{W wielkości miłosierdzia twego wysłuchaj mię,~\red{*} w prawdzie zbawienia twego:}%

\latN{\mred{18.}~Eripe me de luto, ut non in\textbf{fí}gar:~\red{*} líbera me ab iis, qui odérunt me, et de profún\underline{\textit{dis}} \underline{\textit{a}}\textbf{quá}rum.}\polPs{Wyrwij mię z błota, abych nie ulgnął:~\red{*} wybaw mię od tych, którzy mię nienawidzą, i z głębokości wód.}%

\latN{\mred{19.}~Non me demérgat tempéstas aquæ,~† neque absórbeat me pro\textbf{fún}dum:~\red{*} neque úrgeat super me púte\underline{\textit{us}} \underline{\textit{os}} \textbf{su}um.}\polPs{Niech mię nie zatapia nawałność wody, ani mię niech nie pożera głębokość:~\red{*} ani niech nie zawiera nademną studnia wierzchu swego.}%

\latN{\mred{20.}~Exáudi me, Dómine, quóniam benígna est misericórdia \textbf{tu}a:~\red{*} secúndum multitúdinem miseratiónum tuárum ré\underline{\textit{spi}}\underline{\textit{ce}} \textbf{in} me.}\polPs{Wysłuchaj mię, Panie; bo łaskawe jest miłosierdzie twoje:~\red{*} według mnóstwa litości twoich wejrzyj na mię.}%

\latN{\mred{21.}~Et ne avértas fáciem tuam a púero \textbf{tu}o:~\red{*} quóniam tríbulor, velóci\underline{\textit{ter}} \underline{\textit{ex}}\textbf{áu}di me.}\polPs{A nie odwracaj oblicza twego od sługi swego:~\red{*} bom jest w utrapieniu: prędko wysłuchaj mię.}%

\latN{\mred{22.}~Inténde ánimæ meæ, et líbera \textbf{e}am:~\red{*} propter inimícos meos \underline{\textit{é}}\underline{\textit{ri}}\textbf{pe} me.}\polPs{Przybliż się ku duszy mojéj, a wybaw ją:~\red{*} dla nieprzyjaciół moich wyrwij mię.}%

\latN{\mred{23.}~Tu scis impropérium meum, et confusiónem \textbf{me}am,~\red{*} et reverén\underline{\textit{ti}}\underline{\textit{am}} \textbf{me}am.}\polPs{Ty znasz pohańbienie moje i zelżywość moję,~\red{*} i wstyd mój.}%

\latN{\mred{24.}~In conspéctu tuo sunt omnes qui tríbu\textbf{lant} me:~\red{*} impropérium exspectávit cor meum, \underline{\textit{et}} \underline{\textit{mi}}\textbf{sé}riam.}\polPs{Przed oczyma twemi są wszyscy, którzy mię trapią:~\red{*} urągania i nędze czekało serce moje.}%

\latN{\mred{25.}~Et sustínui qui simul contristarétur, et non \textbf{fu}it:~\red{*} et qui consolarétur, et \underline{\textit{non}} \underline{\textit{in}}\textbf{vé}ni.}\polPs{I czekałem, ktoby się społem smęcił, a nie było:~\red{*} i ktoby pocieszył, a nie znalazłem.}%
\EndParallel

\gresetinitiallines{0}\gabcsnippet{(c4) <c><v>\footnotesize{26. }</v></c>Et(j) de(j)dé(j)runt(j) in(j) es(j)cam(j) <b>me</b>(k)am(jr) fel:(j.) <c><v>*</v></c> (:) et(j) in(j) si(j)ti(j) me(j)a(j) po(j)ta(j)vé(j)runt(j) <i><ul>me</i></ul>(h) <i><ul>a</i></ul>(j)<b>cé</b>(k jr)to.(j.) (::)}
\vspace{-8mm}
\tlumaczR{I dali żółć na pokarm mój:~\red{*} a w pragnieniu mojem napawali mię octem.}\linespace%
\BeginParallel
\latN{\mred{27.}~Fiat mensa eórum coram ipsis in \textbf{lá}queum,~\red{*} et in retributiónes, \underline{\textit{et}} \underline{\textit{in}} \textbf{scán}dalum.}\polPs{Niechaj będzie stół ich przed nimi sidłem,~\red{*} i zapłaty, i na upadek.}%

\latN{\mred{28.}~Obscuréntur óculi eórum ne \textbf{ví}deant:~\red{*} et dorsum eórum sem\underline{\textit{per}} \underline{\textit{in}}\textbf{cúr}va.}\polPs{Niech się zaćmią oczy ich, aby nie widzieli:~\red{*} a grzbietu ich zawżdy nachylaj.}%

\latN{\mred{29.}~Effúnde super eos iram \textbf{tu}am:~\red{*} et furor iræ tuæ compre\underline{\textit{hén}}\underline{\textit{dat}} \textbf{e}os.}\polPs{Wyléj na nie gniew twój:~\red{*} i zapalczywość gniewu twego niech je ogarnie.}%

\latN{\mred{30.}~Fiat habitátio eórum de\textbf{sér}ta:~\red{*} et in tabernáculis eórum non sit \underline{\textit{qui}} \underline{\textit{in}}\textbf{há}bitet.}\polPs{Mieszkanie ich niech się stanie puste:~\red{*} a w przybytkach ich niechaj nie będzie, ktoby mieszkał.}%

\latN{\mred{31.}~Quóniam quem tu percussísti, perse\textbf{cú}ti sunt:~\red{*} et super dolórem vúlnerum meórum \underline{\textit{ad}}\underline{\textit{di}}\textbf{dé}runt.}\polPs{Bo, któregoś ty zranił, prześladowali:~\red{*} a na boleści ran moich naddawali.}%

\latN{\mred{32.}~Appóne iniquitátem super iniquitátem e\textbf{ó}rum:~\red{*} et non intrent in justí\underline{\textit{ti}}\underline{\textit{am}} \textbf{tu}am.}\polPs{Przydajże nieprawość ku nieprawości ich:~\red{*} a niech nie wchodzą do sprawiedliwości twojéj.}%

\latN{\mred{33.}~Deleántur de libro vi\textbf{vén}tium:~\red{*} et cum justis \underline{\textit{non}} \underline{\textit{scri}}\textbf{bán}tur.}\polPs{Niechaj będą wymazani z ksiąg żywiących:~\red{*} a niech z sprawiedliwymi nie będą wspisani.}%

\latN{\mred{34.}~Ego sum pauper et \textbf{do}lens:~\red{*} salus tua, De\underline{\textit{us}}, \underline{\textit{su}}\textbf{scé}pit me.}\polPs{Jamci jest ubogi i zbolały:~\red{*} zbawienie twoje, Boże, wspomogło mię.}%

\latN{\mred{35.}~Laudábo nomen Dei cum \textbf{cán}tico:~\red{*} et magnificábo e\underline{\textit{um}} \underline{\textit{in}} \textbf{lau}de:}\polPs{Będę chwalił imię Boże pieśnią:~\red{*} i wielbić go będę chwałą.}%

\latN{\mred{36.}~Et placébit Deo super vítulum no\textbf{vél}lum:~\red{*} córnua producén\underline{\textit{tem}} \underline{\textit{et}} \textbf{ún}gulas.}\polPs{A przyjemniéj będzie Panu, niźli młode cielę:~\red{*} któremu rogi i kopyta podrastają.}%

\latN{\mred{37.}~Vídeant páuperes et læ\textbf{tén}tur:~\red{*} quǽrite Deum, et vivet á\underline{\textit{ni}}\underline{\textit{ma}} \textbf{ves}tra.}\polPs{Niech widzą ubodzy, a niech się weselą:~\red{*} szukajcie Boga, a będzie żyła dusza wasza.}%

\latN{\mred{38.}~Quóniam exaudívit páuperes \textbf{Dó}minus:~\red{*} et vinctos suos \underline{\textit{non}} \underline{\textit{de}}\textbf{spé}xit.}\polPs{Albowiem Pan wysłuchał ubogich:~\red{*} a więźniami swymi nie wzgardził.}%

\latN{\mred{39.}~Laudent illum cæli et \textbf{ter}ra,~\red{*} mare et ómnia reptíli\underline{\textit{a}} \underline{\textit{in}} \textbf{e}is.}\polPs{Niechaj go chwalą niebiosa, i ziemia,~\red{*} morze i wszystek płaz na nich.}%

\latN{\mred{40.}~Quóniam Deus salvam fáciet \textbf{Si}on:~\red{*} et ædificabúntur civi\underline{\textit{tá}}\underline{\textit{tes}} \textbf{Ju}da.}\polPs{Albowiem Bóg zbawi Syon:~\red{*} i będą zbudowane miasta Juda.}%

\latN{\mred{41.}~Et inhabitábunt \textbf{i}bi,~\red{*} et hereditáte ac\underline{\textit{quí}}\underline{\textit{rent}} \textbf{e}am.}\polPs{I będą tam mieszkać,~\red{*} i dostaną jéj dziedzictwem.}%

\latN{\mred{42.}~Et semen servórum ejus possidébit \textbf{e}am:~\red{*} et qui díligunt nomen ejus, habitá\underline{\textit{bunt}} \underline{\textit{in}} \textbf{e}a.}\polPs{I nasienie sług jego otrzymają:~\red{*} a którzy miłują imię jego, będą w niéj mieszkać.}%
\EndParallel
\newpage
\gresetinitiallines{1}\greannotation{1. Ant.}\greannotation{8 c}\gabcsnippet{(c4)ZE(j)lus(j) do(j)mus(j) tu(j_k)æ(ji/jkj') com(h)é(j')dit(i) me,(g.) (;) et(f) op(h)pró(j')bri(i)a(g_[uh:l]h) ex(f)pro(h)brán(j')ti(k)um(i') ti(j)bi(h.) (,) ce(h)ci(g)dé(hi)runt(h') su(g)per(fg~) me.(g.) (::)}
\tlumaczR{Gryzła mnie żarliwość domu twego i urągania urągających tobie na mnie spadły.}
%% -------- 2. Ant. Avertantur retrorsum -----------
\gresetinitiallines{1}\greannotation{2. Ant.}\greannotation{8 c}\gabcsnippet{(c4) A(j)ver(j)tán(j_k)tur(k') re(k)trór(kj~)sum,(k_[hl:1]j_[hl:1]) <c><v>*</v></c>(,) et(j) e(ji)ru(h!iwj)bé(jkj)scant,(g.) (;) qui(j) có(i!kjk)gi(h)tant(f_h) mi(j_h)hi(i) ma(g.)la.(g.) (::)}
\vspace{-8mm}
\tlumaczR{Niech się obrócą wstecz~\red{*} i zawstydzą się, którzy mi chcą zła.}
\subsekcja{Psalm 69}
\gresetinitiallines{0}\gabcsnippet{(c4)<c><v>\footnotesize{1. }</v></c> De(g)us,(h) in(j) ad(j)ju(j)tó(j)ri(j)um(j) me(j)um(j) in(j)<b>tén</b>(k jr)de:(j.) <c><v>*</v></c>(:) Dó(j)mi(j)ne(j) ad(j) ad(j)ju(j)ván(j)dum(j) <i><ul>me</ul></i>(h) <i><ul>fes</ul></i>(j)<b>tí</b>(k jr)na.(j.) ( ::) <c><i>Flex:</i></c>()  tem(j)pés(j)tas(j) a(j)quæ, <c><v>†</v></c>(h. j j j  ::)}
\tlumaczR{Boże, wejrzyj ku wspomożeniu memu:~\red{*} Panie, pośpiesz się ku ratunku memu.}\linespace
\columnratio{0.7}
\setlength{\columnseprule}{0.4pt}
\setlength{\columnsep}{.02\textwidth}
\BeginParallel
\latN{\mred{2.}~Confundántur et revere\textbf{án}tur,~\red{*} qui quærunt á\underline{\textit{ni}}\underline{\textit{mam}} \textbf{me}am.}\polPs{Niechaj będą zawstydzeni i pohańbieni,~\red{*} którzy szukają dusze mojéj.}%

\latN{\mred{3.}~Avertántur retrórsum, et eru\textbf{bés}cant,~\red{*} qui volunt \underline{\textit{mi}}\underline{\textit{hi}} \textbf{ma}la.}\polPs{Niech się obrócą nawstecz, a niech się zawstydzą,~\red{*} którzy mi chcą złego.}%

\latN{\mred{4.}~Avertántur statim erube\textbf{scén}tes,~\red{*} qui dicunt mihi: \underline{\textit{Eu}}\underline{\textit{ge}}, \textbf{eu}ge.}\polPs{Niech się natychmiast obrócą wstydać się,~\red{*} którzy mi mówią: Ehéj, Ehéj.}%

\latN{\mred{5.}~Exsúltent et læténtur in te omnes qui \textbf{quæ}runt te,~\red{*} et dicant semper: Magnificétur Dóminus: qui díligunt salu\underline{\textit{tá}}\underline{\textit{re}} \textbf{tu}um.}\polPs{Niech się weselą, i radują w tobie wszyscy, którzy cię szukają,~\red{*} a niech zawżdy mówią: Niechaj uwielbion będzie Pan: którzy miłują zbawienie twoje.}%

\latN{\mred{6.}~Ego vero egénus, et \textbf{pau}per sum:~\red{*} Deus, \underline{\textit{ád}}\underline{\textit{ju}}\textbf{va} me.}\polPs{A jam jest ubogi i żebrak:~\red{*} Boże, wspomóż mię.}%

\latN{\mred{7.}~Adjútor meus, et liberátor meus \textbf{es} tu:~\red{*} Dómine, \underline{\textit{ne}} \underline{\textit{mo}}\textbf{ré}ris.}\polPs{Pomocnikiem moim, i wybawicielem jesteś ty moim:~\red{*} Panie, nie omieszkajże.}%
%\iffalse %*** %\fi %***
\EndParallel
\segmentspace
\gresetinitiallines{1}\greannotation{2. Ant.}\greannotation{8 c}\gabcsnippet{(c4) A(j)ver(j)tán(j_k)tur(k') re(k)trór(kj~)sum,(k_[hl:1]j_[hl:1]) (,) et(j) e(ji)ru(h!iwj)bé(jkj)scant,(g.) (;) qui(j) có(i!kjk)gi(h)tant(f_h) mi(j_h)hi(i) ma(g.)la.(g.) (::)}
\tlumaczR{Niech się obrócą wstecz i zawstydzą się, którzy mi chcą zła.}
\newpage
\sloppy
\wstawplik{modlitwy_koncowe}

\end{document}
